\paragraph{Exact penalty (E1).} With integer constraint data and bounded binary slack bits, a
penalty weight $\Lam>f^\ast-L(f)$ is sufficient, where $f^\ast$ is the constrained optimum and
$L(f)$ any lower bound. Exhaustive enumeration of the full penalty \QUBO{} (decision plus slack
bits) on 11 instances shows that its ground state equals the MILP optimum on all 11
(Table~\ref{tab:e1}); $\Lsafe$ is 1.81--2.21$\times$ smaller than the range bound.

\paragraph{Penalty weight (E4).} On a mix $8\times8$ + equity instance, $0.01\,\Lsafe$ followed
by repair raises $\Popt$ after post-processing from at most 0.03 to 0.42 (SA) and 0.47 (Tabu).
Raw feasibility collapses below about $0.1\,\Lsafe$. The recipe is hybrid, not certified.

\paragraph{Penalty form (E10, quick run on 2 instances).} An unbalanced penalty at $0.1\,\Lsafe$
needs no slack bits (10 instead of 13 variables) and cuts the dynamic range from 163 to 16.5 with
the ground state still optimal on these instances. It is not certified in general, and the linear
(Lagrangian) form is feasible and optimal only for lucky multipliers.

\paragraph{Quadratisation (E11, quick run on 2 instances).} Rosenberg substitution
\cite{rosenberg1975,boros2002} of a cubic \HUBO{} leaves the slack-penalty final gap unchanged
($9.4\times10^{-5}$ and $1.1\times10^{-5}$) but adds 7 qubits (9 to 16) and lowers the short-anneal
ground-state probability from $3.1\times10^{-3}$ to $2.4\times10^{-5}$. Removing the penalty
altogether (a diagonal-only simulation on $\Fset$) enlarges the gap about 250$\times$.
